\section{Persistence modules}\label{sec:modules}

\begin{definition}[Persistence module]\label{def:pm}
Let $(P,\le)$ be a poset. A \emph{$P$-indexed persistence module} is a functor $M\colon P\to\Vec$.
We write $M_t$ for $M(t)$ and $\varphi_M(s,t)\colon M_s\to M_t$ for the structure maps, following~\cite[Definition~2.1]{chazal2016structure}.
Such an $M$ is \emph{pointwise finite-dimensional} \textup{(}pfd\textup{)} if every $M_t$ is finite-dimensional.
\end{definition}

\begin{definition}[Interleaving distance]\label{def:interleave}
For $\delta\ge 0$ define
\begin{align}
d_I(M,N) &:= \inf\{\delta\ge0: M,N \text{ are } \delta\text{-interleaved}\},\label{eq:di}\\
\rho(M) &= \sup_t \dim M_t. \notag
\end{align}
\end{definition}

\begin{convention}
Throughout, $\Hloc{q}{\mathfrak m}(-)$ denotes local cohomology with support in $\mathfrak m$.
\end{convention}

\begin{example}
Let $\Delta$ be the path $1-2-3$. Then $\link_\Delta(2)=\{\varnothing,\{1\},\{3\}\}$.
\end{example}
